\clearpage
\fancyhead[L]{\small BURCKHARDT. TABLES DE LA LUNE.}
\SourceLeaf{558}
\begin{center}\small Göttingische gelehrte Anzeigen. 39. Stück. Seiten 388..392. 8. März 1813.\\[3pt]\rule{.72\linewidth}{.3pt}\end{center}

\emph{Tables astronomiques publiées par le bureau des longitudes de France. Tables de la lune, par} M. \textsc{Burckhardt}, \emph{membre de l'institut impérial, du bureau des longitudes de France et de plusieurs autres sociétés savantes. Décembre} 1812. Paris. Bei Madame Courcier. 88 Seiten in Quart.

Die Berechnung der Mondpositionen hat durch die Tafeln von \textsc{Bürg} einen so hohen Grad an Genauigkeit erreicht, dass es ein gewagtes Unternehmen scheint, diese berühmten Tafeln noch übertreffen zu wollen. Die größtmögliche Vollkommenheit der Mondtafeln ist jedoch in vieler Hinsicht so wichtig, dass man einem so geschickten Astronomen wie dem Verfasser der vorliegenden Tafeln für seine Bemühungen um weitere Verbesserungen den größten Dank schuldet. Arbeiten dieser Art verdienen umso mehr Anerkennung, als ihnen statt einiger Monate ganze \emph{Jahre} geopfert werden müssen und sie der Natur der Sache nach nicht mehr durch \emph{glänzende} Erfolge belohnt werden können.

Der Verfasser verfolgte mit seinem Unternehmen zwei Ziele. Zunächst wollte er den Tafeln eine etwas veränderte, bequemere und einfachere \emph{Form} geben. Dazu begnügte er sich nicht damit, die Elemente der Mondtafeln von \textsc{Bürg} umzugießen, sondern legte seinen eigenen Tafeln eine selbst vorgenommene neue Auswertung von mehr als viertausend Beobachtungen zugrunde. Diese Tafeln sind daher als wirklich alleiniges Eigentum von Herrn \textsc{Burckhardt} anzusehen. Er durfte zu Recht hoffen, dass die neuen Tafeln so auch in der \emph{Genauigkeit} einen gewissen Vorsprung gewinnen würden. Das bestätigt tatsächlich der Vergleich von dreihundert Beobachtungen, den das französische Büro der Längengrade sowohl mit den Tafeln von \textsc{Burckhardt} als auch mit denen von \textsc{Bürg} vornehmen ließ. Für 167 in Greenwich und an der kaiserlichen Sternwarte in Paris beobachtete Längen betrug die Summe der Fehlerquadrate nach den Tafeln von \textsc{Burckhardt} $= 4602''$, nach denen von \textsc{Bürg} dagegen $= 7083''$. Für 137 andere in Paris an der kaiserlichen Sternwarte und an der Militärschule beobachtete Längen betrug die Summe der Fehlerquadrate nach den Tafeln von \textsc{Burckhardt} $4182''$, nach denen von \textsc{Bürg} $6439''$. Auch die Breiten stimmten, wie der Verfasser versichert, mit seinen eigenen Tafeln besser überein als mit denen von \textsc{Bürg}. Die Größe der Abweichung ist hier jedoch nicht angegeben. Wir hätten die Ergebnisse all dieser Vergleiche gern einzeln abgedruckt gesehen. Dadurch wäre die hohe Vollkommenheit der Mondtafeln, die inzwischen erreicht ist, noch anschaulicher geworden. Außerdem wäre die künftige Prüfung, ob die eine oder andere zusätzliche Gleichung die Tafeln noch verbessert, ungemein erleichtert worden. Überhaupt hätten wir uns in der Einleitung zu diesen Tafeln etwas mehr Ausführlichkeit gewünscht. Zwar sind wir keineswegs für die weitschweifigen Anleitungen, in denen manche Verfasser von Tafeln allgemein Bekanntes bis zum Überdruss wiederholen. Doch scheint es uns ein berechtigter Wunsch, die Besonderheiten neuer Tafeln knapp, aber vollständig und ausdrücklich angegeben zu finden, auch wenn sie sich durch Analyse der Tafeln selbst leicht ermitteln lassen.

Der wichtigste Unterschied zwischen der Form dieser Tafeln und der seit \textsc{Tobias Meyer} von \textsc{Mason}, \textsc{Triesnecker} und \textsc{Bürg} beibehaltenen Form besteht darin, dass ihnen statt der wahren die mittleren Sonnenlängen zugrunde liegen. Auch die Länge des Knotens und das Perigäum werden hier nicht erst durch eine Gleichung korrigiert, die von der Sonnenanomalie abhängt. Die Evektion ist dagegen von den übrigen kleinen Gleichungen getrennt. Die Zahl der kleinen Gleichungen, deren Argumente \SourceLeaf{559}sich alle unabhängig voneinander leicht und unmittelbar aus der Tafel entnehmen lassen, beträgt jetzt 32, einschließlich der Nutation und zweier Störungsgleichungen von Venus und Jupiter. Damit wird das mittlere Argument der Evektion korrigiert. Die Summe dieser 32 Gleichungen und der Evektion korrigiert die Anomalie. Die Summe der 32 Gleichungen, der Evektion und der Mittelpunktsgleichung korrigiert das Argument der Variation. Die Variation ergibt zusammen mit den vorigen 34 Gleichungen und der mittleren Länge die wahre Länge in der Bahn, die dann noch auf die Ekliptik reduziert wird. Bei den Argumenten der Evektion, der Mittelpunktsgleichung und der Variation sowie bei der mittleren Länge ist zugleich die Säkulargleichung zu berücksichtigen. Mit Letzterer ist die bekanntlich empirisch bestimmte kleine Ungleichheit langer Periode in den Epochen für das 19. Jahrhundert zusammengefasst. Für andere Jahrhunderte sind beide nach einer am Schluss hinzugefügten besonderen Tafel getrennt zu berechnen. Für diese kleine Gleichung hält es \textsc{Laplace} jetzt für am wahrscheinlichsten, dass sie dem Kosinus der Summe aus der doppelten Länge des Mondknotens und der Länge des Perigäums proportional ist. Die Periode beträgt demnach 175 Jahre, der Koeffizient wird $= 12''5$ gesetzt, und die Gleichung hat ihren Ursprung in einer angenommenen Ungleichheit zwischen der nördlichen und der südlichen Erdhalbkugel. Es ist zu wünschen, dass die Theorie in diesem wichtigen Punkt noch weiter vervollkommnet wird. In den Mondtafeln von \textsc{Bürg} nahm man an, diese Gleichung sei dem Sinus jenes um die dreifache Länge des Sonnenperigäums verminderten Arguments proportional. Ihr Koeffizient wurde $= 14''$ angenommen. Bei der Horizontalparallaxe hat sich Herr \textsc{Burckhardt} ganz an die Theorie von \textsc{Laplace} gehalten. Das Verhältnis dieser Parallaxe zum horizontalen Halbmesser des Mondes beruht auf der bei Vollmond beobachteten Dauer der Meridiandurchgänge. --- Am Ende des Werks sind noch zwei von Herrn \textsc{Burckhardt} neu entwickelte Formeln angefügt, die sich mit Vorteil zur Berechnung der Neu- und Vollmonde verwenden lassen.
